% ================================================================
% chapters/ch3_methodology.tex
% Chapter 3 — Materials and Methods
% Sections 3.1 to 3.4 (this file)
% ================================================================

\chapter{Materials and Methods}
\label{ch:methodology}

This chapter presents the complete theoretical framework and
implementation of the metadata de-anonymization risk model.
It is structured in two parts. The first part --- Sections
\ref{sec:system_model} through \ref{sec:mitigation} ---
derives the mathematical model from first principles. The
second part --- Sections \ref{sec:fabric_context} through
\ref{sec:implementation} --- describes the system design and
Python implementation through which the model is applied to
Nepal's domestic CBDC context.

% ================================================================
\section{System Model}
\label{sec:system_model}
% ================================================================

\subsection{Population and Users}

Let $U$ denote the set of all users in the CBDC system,
with $|U| = N$ where $N$ is the total population size. Each
user $u \in U$ represents a distinct real-world individual
who holds a CBDC wallet and participates in the domestic
payment system.

In this model, the population is treated as a fixed snapshot
at a given point in time. This reflects the scenario where
an adversary obtains the complete transaction ledger and
attempts to identify users from the observable metadata.
The scalability analysis in Chapter~\ref{ch:results}
examines how the risk profile changes as $N$ varies from
a small village pilot to a national deployment.

\subsection{Attributes and Value Domains}

Each transaction in the CBDC system is associated with a
set of $m$ observable metadata attributes. Formally, each
attribute $a_i$ is a function that maps every user in $U$
to a value in a domain $V_i$:

\begin{equation}
a_i : U \rightarrow V_i \quad \text{for } i = 1, 2, \ldots, m
\label{eq:attribute_function}
\end{equation}

The value domain $V_i$ is the complete set of all possible
values that attribute $a_i$ can take. The nature of $V_i$
depends on the attribute:

\begin{itemize}
    \item For \textit{amount}, $V_i \subset \mathbb{R}^+$,
    the positive reals. Under Plain Hash Token, $V_i$ is
    effectively continuous. Under ZKP Token, $V_i$ is a
    finite set of five denomination bands.

    \item For \textit{timestamp}, $V_i \subset \mathbb{Z}$,
    the integers. The granularity --- from fine-grained
    Unix-style integers to day-of-week only --- depends on
    the cryptographic scheme.

    \item For \textit{merchant}, \textit{payment channel},
    and \textit{transaction zone}, $V_i$ is a finite
    categorical set whose cardinality varies by scheme.

    \item For \textit{wallet}, $V_i$ is a finite set of
    address identifiers whose size depends on the scheme's
    wallet generation mechanism.

    \item For \textit{frequency}, $V_i \subset \mathbb{Z}^+$,
    the positive integers, bounded by the scheme's
    transaction rate parameters.
\end{itemize}

The value domain is determined at three levels. Reality sets
the outer boundary: amounts cannot be negative, zones cannot
refer to regions outside Nepal. The CBDC system design sets
the precision and granularity: the chaincode records amounts
to two decimal places under Plain Hash Token or to the
nearest NPR 100 under Ring Signature. The cryptographic
scheme then transforms what is actually observable: ZKP range
proofs reduce the observable amount to a coarse band
regardless of the underlying transaction value.

\subsection{The Attribute Matrix}

The complete metadata record of the CBDC system can be
represented as a matrix $\mathbf{A}$ of dimension $N \times m$,
where entry $\mathbf{A}_{u,i} = a_i(u)$ is the value of
attribute $i$ for user $u$. Each row corresponds to a user's
transaction fingerprint. The adversary's goal is to identify
which row corresponds to a target individual --- that is,
to invert the mapping from observable metadata to real
identity.

% ================================================================
\section{Adversary Model}
\label{sec:adversary_model}
% ================================================================

\subsection{Adversary Type}

The adversary is modelled as a \textit{passive,
honest-but-curious observer}. This is a standard adversary
model in cryptographic privacy literature. The adversary:

\begin{itemize}
    \item \textbf{Observes} all $m$ metadata fields for
    every transaction on the ledger. They have complete
    read access to the attribute matrix $\mathbf{A}$.

    \item \textbf{Does not modify} transactions, inject
    false records, or perform active attacks on the
    cryptographic primitives.

    \item \textbf{Has unlimited computation}. There is no
    computational hardness assumption in this model. The
    threat is purely information-theoretic: what can be
    inferred from the structure of the observable data,
    regardless of computational cost.

    \item \textbf{Has one goal}: given a transaction record
    --- a row of $\mathbf{A}$ --- identify which user
    $u \in U$ produced it.
\end{itemize}

\subsection{Why Unlimited Computation}

The unlimited computation assumption is deliberate and
important. Computational hardness assumptions --- the basis
of most cryptographic security proofs --- are time-dependent.
What is computationally infeasible today may become feasible
with advances in hardware, algorithms, or quantum computing.
By assuming unlimited computation, this model produces
risk bounds that are permanent: if the information is
structurally present in the metadata, the adversary can
exploit it regardless of when or with what hardware they
attempt to do so.

This is an information-theoretic threat model, asking not
``can the adversary compute this?'' but ``is the information
present in the data at all?'' The answer to the latter
question does not change over time.

\subsection{Concrete Adversary Instantiations}

The passive honest-but-curious adversary model covers three
realistic threat scenarios in Nepal's CBDC context.

\textbf{The compromised insider.} A Nepal Rastra Bank
employee with read access to the transaction database. They
have not broken any cryptographic system --- they have
legitimate access --- but they attempt to use transaction
patterns to build behavioural profiles of specific users.
This is not a theoretical risk: central bank data breaches
have occurred in Bangladesh, Ecuador, and multiple other
jurisdictions.

\textbf{The passive network participant.} Any organisation
that is a member of the Hyperledger Fabric channel sees all
transactions on that channel. A large merchant network,
a payment processor, or a correspondent bank that is a
channel participant has full read access to the ledger
without any attack.

\textbf{The merchant coalition.} Multiple merchants
independently see their own transactions with users. By
sharing transaction data --- amounts, times, wallet
addresses --- they can construct a combined view of user
behaviour across platforms. Nepal's dominant payment
platforms eSewa and Khalti already hold significant
transaction data on the same user populations that would
use a CBDC.

\subsection{Scope of the Adversary Model}

This model does not cover active adversaries who inject
false transactions, perform Sybil attacks, or attempt to
break cryptographic primitives by brute force or algorithm.
Those threat categories are outside the scope of this
thesis and represent a separate research problem. The
passive metadata analysis threat is chosen because it is
the realistic and immediate risk for a CBDC deployment ---
no sophisticated attack is required, only the ability to
read and analyse what is already publicly visible on the
ledger.

% ================================================================
\section{Single-Attribute Uniqueness Score \texorpdfstring{$s(a_i)$}{s(ai)}}
\label{sec:s_attr}
% ================================================================

\subsection{Definition}

The single-attribute uniqueness score $s(a_i)$ measures the
fraction of users in the population who are uniquely
identifiable from attribute $a_i$ alone.

A user $u$ is \textit{unique} on attribute $a_i$ if no
other user in $U$ shares the same attribute value:

\begin{equation}
\mathbb{I}_i(u) =
\begin{cases}
1 & \text{if } |\{u' \in U : a_i(u') = a_i(u)\}| = 1 \\
0 & \text{otherwise}
\end{cases}
\label{eq:indicator}
\end{equation}

The indicator function $\mathbb{I}_i(u)$ equals 1 if and
only if user $u$ is the sole occupant of their attribute
value --- no other user shares it. The uniqueness score is
then the population-weighted average of this indicator:

\begin{equation}
s(a_i) = \frac{1}{N} \sum_{u \in U} \mathbb{I}_i(u)
\label{eq:s_single}
\end{equation}

\subsection{Bounds and Interpretation}

\begin{proposition}
$0 \leq s(a_i) \leq 1$ for any attribute $a_i$ and any
population $U$.
\end{proposition}

\begin{proof}
The indicator $\mathbb{I}_i(u) \in \{0, 1\}$ for every
$u \in U$. Therefore $\sum_{u \in U} \mathbb{I}_i(u)
\in [0, N]$. Dividing by $N$ gives $s(a_i) \in [0, 1]$.
\end{proof}

The two extremes have precise interpretations:

\begin{itemize}
    \item $s(a_i) = 0$ means no user is unique on attribute
    $a_i$. Every user shares their attribute value with at
    least one other user. The adversary gains no identifying
    information from observing this attribute. This is the
    privacy ideal for any single attribute.

    \item $s(a_i) = 1$ means every user is unique on
    attribute $a_i$. Every user has a distinct value that
    appears in exactly one record. Observing this attribute
    alone is sufficient to identify the user completely.
    This is the worst case for any single attribute.
\end{itemize}

Values between 0 and 1 represent partial risk. $s(a_i) =
0.317$ means 31.7\% of users are uniquely identifiable
from attribute $a_i$ alone; the remaining 68.3\% share
their attribute value with at least one other user and
are protected from single-attribute identification.

\subsection{Connection to k-Anonymity}

The single-attribute uniqueness score connects formally
to the k-anonymity framework introduced by Sweeney
\citeyear{sweeney2002}. A user satisfies k-anonymity
on attribute $a_i$ with $k \geq 2$ if they share their
attribute value with at least $k-1$ other users. A user
who is unique --- $k=1$ --- has zero k-anonymity protection
on that attribute.

The uniqueness score $s(a_i)$ is precisely the
\textit{k-anonymity failure rate} for attribute $a_i$:
the fraction of users for whom k-anonymity has completely
collapsed to $k=1$. This connection provides theoretical
grounding for the metric: maximising k-anonymity across
the population is equivalent to minimising $s(a_i)$
toward zero.

\subsection{Worked Example}

Consider a population of $N = 5$ users and the attribute
\textit{transaction zone}:

\begin{table}[H]
\centering
\caption{Example: single-attribute uniqueness on zone}
\label{tab:s_example}
\begin{tabular}{lll}
\toprule
\textbf{User} & \textbf{Zone} & \textbf{Unique?} \\
\midrule
Alice  & Kathmandu valley & No --- shared with Carol, Eve \\
Bob    & Hill rural       & Yes --- sole occupant \\
Carol  & Kathmandu valley & No --- shared with Alice, Eve \\
David  & Terai urban      & Yes --- sole occupant \\
Eve    & Kathmandu valley & No --- shared with Alice, Carol \\
\bottomrule
\end{tabular}
\end{table}

$\sum \mathbb{I}_i(u) = 0 + 1 + 0 + 1 + 0 = 2$

$s(\text{zone}) = 2/5 = 0.400$

Bob and David are uniquely identifiable from zone alone.
Alice, Carol, and Eve are protected by sharing the
Kathmandu valley value. $s = 0.400$ means 40\% of this
population is exposed by this attribute.

% ================================================================
\section{Pairwise Interaction Coefficient \texorpdfstring{$\lambda_{ij}$}{lambda\_ij}}
\label{sec:lambda}
% ================================================================

\subsection{Motivation}

Single-attribute uniqueness captures only part of the
privacy risk. Consider a user who is not unique on zone
alone and not unique on payment channel alone. If the
combination of their zone and payment channel values is
unique --- shared by no other user --- they are fully
identifiable from the pair of attributes even though
neither attribute exposes them individually. This is the
attribute combination effect demonstrated empirically by
Narayanan and Shmatikoff \citeyear{narayanan2008} in the
Netflix de-anonymization study. The pairwise interaction
coefficient $\lambda_{ij}$ formally captures this
amplification.

\subsection{Joint Uniqueness}

Define the \textit{joint uniqueness score} $s(a_i, a_j)$
as the fraction of users who are unique on the
\textit{combination} of attributes $a_i$ and $a_j$
simultaneously:

\begin{equation}
s(a_i, a_j) = \frac{1}{N}
\sum_{u \in U} \mathbb{I}_{ij}(u)
\label{eq:s_joint}
\end{equation}

where $\mathbb{I}_{ij}(u) = 1$ if and only if no other
user $u' \in U$ shares both $a_i(u') = a_i(u)$ and
$a_j(u') = a_j(u)$ simultaneously. A user must match on
both attributes to be considered non-unique.

\subsection{Definition of $\lambda_{ij}$}

If attributes $a_i$ and $a_j$ were statistically
independent, the expected joint uniqueness under
independence would be $s(a_i) \times s(a_j)$ --- the
product of the individual uniqueness scores. The pairwise
interaction coefficient $\lambda_{ij}$ measures how much
the actual joint uniqueness deviates from this
independence baseline:

\begin{equation}
\lambda_{ij} = \frac{s(a_i, a_j)}{s(a_i) \times s(a_j)}
\label{eq:lambda}
\end{equation}

\subsection{Interpretation}

The three cases of $\lambda_{ij}$ have clear interpretations:

\begin{itemize}
    \item $\lambda_{ij} = 1$: the attributes contribute
    \textit{independently}. The combination identifies
    exactly as many users as independence predicts. No
    amplification.

    \item $\lambda_{ij} > 1$: \textit{amplified risk}.
    The combination is more identifying than independence
    predicts. Users who are safe on each attribute
    individually become exposed by the combination. This
    is the Netflix effect.

    \item $\lambda_{ij} < 1$: \textit{suppressed risk}.
    The attributes are redundant --- they carry overlapping
    information. The combination identifies fewer additional
    users than independence predicts.
\end{itemize}

\subsection{Special Case: Zero Denominator}

When $s(a_i) = 0$ or $s(a_j) = 0$, the denominator of
Equation~\ref{eq:lambda} is zero and $\lambda_{ij}$ is
undefined. By convention, $\lambda_{ij} = 0$ in this case.

This convention is justified by examining the contribution
of the pair to the risk score $\Rraw$:

\begin{equation}
\text{pair contribution} =
\lambda_{ij} \times s(a_i) \times s(a_j)
\end{equation}

If either $s(a_i) = 0$ or $s(a_j) = 0$, this product
equals zero regardless of $\lambda_{ij}$. Setting
$\lambda_{ij} = 0$ produces the correct result --- zero
pair contribution --- without division by zero. In the
Python implementation this is handled as:

\begin{lstlisting}
def lam(si, sj, sij):
    d = si * sj
    return 0.0 if d < 1e-12 else sij / d
\end{lstlisting}

The threshold $10^{-12}$ rather than exact zero accounts
for floating point precision.

\subsection{Symmetry}

\begin{proposition}
$\lambda_{ij} = \lambda_{ji}$ for all pairs $(i, j)$.
\end{proposition}

\begin{proof}
The joint uniqueness score is symmetric by definition:
$s(a_i, a_j) = s(a_j, a_i)$, since the condition
``no other user shares both values'' is invariant to
the order in which attributes are named. The denominator
$s(a_i) \times s(a_j) = s(a_j) \times s(a_i)$ by
commutativity of multiplication. Therefore
$\lambda_{ij} = \lambda_{ji}$.
\end{proof}

This symmetry means the interaction matrix
$\Lambda = [\lambda_{ij}]$ is symmetric. In the
$\lambda_{ij}$ heatmaps in Chapter~\ref{ch:results},
the matrix is visually symmetric across the diagonal.

\subsection{Worked Example}

Extending the five-user example, add a second attribute
\textit{payment channel}:

\begin{table}[H]
\centering
\caption{Example: joint uniqueness and $\lambda_{ij}$}
\label{tab:lambda_example}
\begin{tabular}{lllll}
\toprule
\textbf{User} & \textbf{Zone} & \textbf{Channel} &
\textbf{Pair} & \textbf{Joint unique?} \\
\midrule
Alice  & Kathmandu & Mobile & (KTM, Mob) & No --- Eve same \\
Bob    & Hill rural & Mobile & (HR, Mob)  & Yes \\
Carol  & Kathmandu & QR     & (KTM, QR)  & Yes \\
David  & Terai urb & Mobile & (TU, Mob)  & Yes \\
Eve    & Kathmandu & Mobile & (KTM, Mob) & No --- Alice same \\
\bottomrule
\end{tabular}
\end{table}

From the earlier example: $s(\text{zone}) = 0.400$,
$s(\text{channel}) = 1/5 = 0.200$ (Carol is the only
QR user).

$s(\text{zone, channel}) = 3/5 = 0.600$
(Bob, Carol, David are jointly unique).

\begin{equation*}
\lambda_{\text{zone,channel}} =
\frac{0.600}{0.400 \times 0.200} = \frac{0.600}{0.080}
= 7.5
\end{equation*}

The combination is 7.5 times more identifying than
independence predicts. Carol was protected on zone alone
(shared Kathmandu) but exposed by the combination
(sole Kathmandu QR user). This 7.5$\times$ amplification
is the core of the pairwise interaction effect.

% ================================================================
\section{Raw Risk Score \texorpdfstring{$R_{\mathrm{raw}}$}{R\_raw}}
\label{sec:r_raw}
% ================================================================

\subsection{From Two Attributes to m Attributes}

The two-attribute formulation combines individual
uniqueness and pairwise interaction into a single
raw risk accumulator:

\begin{equation}
R = s(a_1) + s(a_2) + \lambda_{12} \times s(a_1) \times s(a_2)
\label{eq:r_two}
\end{equation}

This expression has three terms: the individual exposure
of attribute 1, the individual exposure of attribute 2,
and the amplified joint exposure from their combination.
Generalising to $m$ attributes requires extending both
sums. For $m$ attributes there are $m$ individual terms
and $\binom{m}{2} = m(m-1)/2$ distinct pairs.

\subsection{General Definition}

The raw risk score for $m$ attributes is:

\begin{equation}
\Rraw = \sum_{i=1}^{m} s(a_i)
\;+\;
\sum_{i=1}^{m} \sum_{j > i} \lambda_{ij} \times s(a_i) \times s(a_j)
\label{eq:r_raw}
\end{equation}

The first sum accumulates individual attribute leakage
across all $m$ attributes. The second sum accumulates
pairwise amplification effects across all $\binom{m}{2}$
distinct pairs. The condition $j > i$ ensures each pair
$(i,j)$ is counted exactly once --- the pair
$(\text{amount}, \text{wallet})$ is counted when $i =
\text{amount}$ and $j = \text{wallet}$ but not again
when $i = \text{wallet}$ and $j = \text{amount}$, since
$\lambda_{ij} = \lambda_{ji}$ by symmetry.

\subsection{Structure of the Pair Terms}

Each pair term $\lambda_{ij} \times s(a_i) \times s(a_j)$
equals $s(a_i, a_j)$ by the definition of $\lambda_{ij}$:

\begin{equation}
\lambda_{ij} \times s(a_i) \times s(a_j)
= \frac{s(a_i, a_j)}{s(a_i) \times s(a_j)}
\times s(a_i) \times s(a_j)
= s(a_i, a_j)
\label{eq:pair_term}
\end{equation}

This means $\Rraw$ can equivalently be read as the sum of
all individual uniqueness scores plus the sum of all joint
uniqueness scores across every distinct attribute pair.
The $\lambda_{ij}$ formulation is retained because it
makes the amplification beyond independence explicit and
measurable.

\subsection{Counting the Pairs}

For this thesis with $m = 7$ attributes, the number of
distinct pairs is:

\begin{equation}
\binom{7}{2} = \frac{7 \times 6}{2} = 21
\end{equation}

The seven individual terms plus the 21 pair terms give
a total of 28 terms in $\Rraw$. This figure is not
coincidental --- it equals $R_{\mathrm{max}}(7) = 28$,
derived in the following section.

\subsection{Why Most Terms Vanish in Practice}

In the results of this thesis, most pair terms contribute
approximately zero to $\Rraw$. This occurs because the
Nepal-specific attributes --- payment channel and
transaction zone --- have $s(a_i) \approx 0$ due to
their skewed distributions. Any pair term involving an
attribute with $s(a_i) = 0$ contributes
$\lambda_{ij} \times 0 \times s(a_j) = 0$ regardless
of the value of $\lambda_{ij}$. For Plain Hash Token,
where $s(\text{wallet}) = 1.000$ and
$s(\text{amount}) = 0.988$, the dominant contribution
comes from the (wallet, amount) pair, the (wallet,
timestamp) pair, and the individual terms for wallet
and amount. The remaining 19 pairs contribute little.

\subsection{Python Implementation}

The \texttt{R\_raw} function in the implementation
directly maps to Equation~\ref{eq:r_raw}:

\begin{lstlisting}
def R_raw(s, lmat, attrs):
    R = sum(s[a] for a in attrs)
    for ai, aj in combinations(attrs, 2):
        R += lmat[ai].get(aj, 0.0) * s[ai] * s[aj]
    return R
\end{lstlisting}

The \texttt{combinations(attrs, 2)} call from Python's
\texttt{itertools} library automatically generates all
$\binom{m}{2}$ pairs with $j > i$, exactly matching the
double sum in Equation~\ref{eq:r_raw}.

% ================================================================
\section{Theoretical Maximum \texorpdfstring{$R_{\mathrm{max}}$}{R\_max}}
\label{sec:r_max}
% ================================================================

\subsection{Motivation}

$\Rraw$ as defined in Equation~\ref{eq:r_raw} is an
un-normalised quantity. Its value depends on $m$: with
more attributes, more terms accumulate and $\Rraw$ grows.
A value of $\Rraw = 8.886$ for Plain Hash Token with
$m = 7$ is meaningless without knowing the scale.
To interpret it, we need to know the maximum value
$\Rraw$ could possibly reach.

$R_{\mathrm{max}}$ is the \textit{theoretical upper bound}
on $\Rraw$ --- the value it achieves under the
worst-case configuration of all attributes simultaneously.
Deriving $R_{\mathrm{max}}$ from first principles is a
core contribution of this thesis.

\subsection{Worst-Case Conditions}

$\Rraw$ is maximised when every term in
Equation~\ref{eq:r_raw} simultaneously reaches its
maximum value. Two conditions must hold jointly:

\textbf{Condition 1:} $s(a_i) = 1$ for all $i = 1,
\ldots, m$. Every user is unique on every attribute
individually. This is the maximum individual exposure
possible.

\textbf{Condition 2:} $\lambda_{ij} = 1$ for all pairs
$(i, j)$. All attributes interact at the independence
level.

Under these conditions each individual term equals:
$s(a_i) = 1$.

Each pair term equals:
$\lambda_{ij} \times s(a_i) \times s(a_j)
= 1 \times 1 \times 1 = 1$.

\subsection{Why $\lambda_{ij} = 1$ Gives the Maximum}

This may appear counterintuitive --- surely $\lambda_{ij}
> 1$ would give larger pair terms? The answer is no,
because of the bound on $s(a_i, a_j)$.

From Equation~\ref{eq:pair_term}, each pair term equals
$s(a_i, a_j)$, which is a fraction of users and
therefore bounded above by 1. When $s(a_i) = s(a_j) = 1$,
setting $\lambda_{ij} = 1$ gives $s(a_i, a_j) = 1$,
which is already the maximum the pair term can achieve.
No value of $\lambda_{ij} > 1$ can push the pair term
above 1 because $s(a_i, a_j) \leq 1$ always. The
bound of 1 per term is tight at $\lambda_{ij} = 1$
when both individual scores are also 1.

\subsection{Derivation}

Under the worst-case conditions:

\begin{align}
R_{\mathrm{max}}(m)
&= \underbrace{\sum_{i=1}^{m} 1}_{m \text{ terms}}
+ \underbrace{\sum_{i=1}^{m} \sum_{j > i} 1}_{\binom{m}{2}
\text{ terms}} \nonumber \\[6pt]
&= m + \binom{m}{2} \nonumber \\[6pt]
&= m + \frac{m(m-1)}{2}
\label{eq:r_max}
\end{align}

\subsection{Values for Relevant $m$}

\begin{table}[H]
\centering
\caption{$R_{\mathrm{max}}(m)$ for $m = 1$ to $m = 7$}
\label{tab:r_max}
\begin{tabular}{cccc}
\toprule
$m$ & Individual terms & Pair terms $\binom{m}{2}$ &
$R_{\mathrm{max}}$ \\
\midrule
1 & 1 & 0  & 1  \\
2 & 2 & 1  & 3  \\
3 & 3 & 3  & 6  \\
4 & 4 & 6  & 10 \\
5 & 5 & 10 & 15 \\
6 & 6 & 15 & 21 \\
7 & 7 & 21 & \textbf{28} \\
\bottomrule
\end{tabular}
\end{table}

For this thesis with $m = 7$:

\begin{equation}
R_{\mathrm{max}}(7) = 7 + \frac{7 \times 6}{2} = 7 + 21 = 28
\end{equation}

\subsection{Quadratic Growth and Its Implication}

$R_{\mathrm{max}}$ grows quadratically in $m$. Each new
attribute added increases $R_{\mathrm{max}}$ by $m$ ---
adding the $m$-th attribute introduces one new individual
term and $m-1$ new pair terms, for a total increase of
$m$. This quadratic growth is the formal explanation for
the counterintuitive finding in the m-sensitivity analysis:
adding Nepal's low-cardinality attributes decreases
$R_{\mathrm{norm}}$ because $R_{\mathrm{max}}$ grows by
6 and then 7, while $\Rraw$ increases by near-zero.

\subsection{Proof that $\Rraw \leq R_{\mathrm{max}}$}

\begin{theorem}
For any valid population $U$, any attribute set of size
$m$, and any cryptographic scheme:
$\Rraw \leq R_{\mathrm{max}}(m)$.
\end{theorem}

\begin{proof}
It suffices to show each term in $\Rraw$ is bounded by
its maximum-case value of 1.

For individual terms: $s(a_i) \leq 1$ for all $i$,
since $s(a_i)$ is a fraction of users.

For pair terms: from Equation~\ref{eq:pair_term},
each pair term equals $s(a_i, a_j)$. Since
$s(a_i, a_j)$ is a fraction of users,
$s(a_i, a_j) \leq 1$ for all pairs.

Therefore each of the $m$ individual terms contributes
at most 1, and each of the $\binom{m}{2}$ pair terms
contributes at most 1. Summing all maximum contributions:
$\Rraw \leq m + \binom{m}{2} = R_{\mathrm{max}}(m)$.
\end{proof}

This theorem guarantees that the normalisation in the
following section always produces a score in $[0, 1]$.

% ================================================================
\section{Normalised Risk Score \texorpdfstring{$R_{\mathrm{norm}}$}{R\_norm}}
\label{sec:r_norm}
% ================================================================

\subsection{Definition}

The normalised risk score $R_{\mathrm{norm}}$ maps
$\Rraw$ into the unit interval by dividing by the
proven theoretical maximum:

\begin{equation}
R_{\mathrm{norm}} = \frac{\Rraw}{R_{\mathrm{max}}(m)}
\label{eq:r_norm}
\end{equation}

\subsection{Properties}

\begin{proposition}
$R_{\mathrm{norm}}$ satisfies the following properties
for any $m$, $N$, and cryptographic scheme.
\end{proposition}

\textbf{Property 1 --- Boundedness:}
$R_{\mathrm{norm}} \in [0, 1]$.

This follows directly from the theorem in
Section~\ref{sec:r_max}: $0 \leq \Rraw
\leq R_{\mathrm{max}}$, and dividing by $R_{\mathrm{max}}
> 0$ preserves the bounds.

\textbf{Property 2 --- Zero case:}
$R_{\mathrm{norm}} = 0$ if and only if $s(a_i) = 0$
for all $i = 1, \ldots, m$.

When every $s(a_i) = 0$, every pair term
$\lambda_{ij} \times s(a_i) \times s(a_j) = 0$,
so $\Rraw = 0$ and $R_{\mathrm{norm}} = 0$.
This represents the privacy ideal: no user is
uniquely identifiable on any attribute or any
combination of attributes.

\textbf{Property 3 --- Maximum case:}
$R_{\mathrm{norm}} = 1$ if and only if
$\Rraw = R_{\mathrm{max}}$, achieved when $s(a_i) = 1$
for all $i$ and $\lambda_{ij} = 1$ for all pairs.
This represents maximum possible risk: every user is
uniquely identifiable on every attribute and every
combination.

\textbf{Property 4 --- Cross-$m$ comparability:}
$R_{\mathrm{norm}}$ is directly comparable across
different values of $m$, because both the numerator
and denominator scale with $m$ through the same
structure. A model with $m = 5$ attributes and
a model with $m = 7$ attributes can be compared
on $R_{\mathrm{norm}}$ without adjustment. This
property makes the m-sensitivity analysis in
Chapter~\ref{ch:results} formally valid.

\textbf{Property 5 --- Monotonicity:}
$R_{\mathrm{norm}}$ is non-decreasing in each
$s(a_i)$ and in each $\lambda_{ij}$, holding all
other quantities fixed. Increasing any individual
uniqueness score or any interaction coefficient
increases $\Rraw$ and therefore $R_{\mathrm{norm}}$.
Privacy never improves when uniqueness increases.

\subsection{Worked Example with Normalisation}

Using the original worked example from the thesis
model specification with $m = 2$ attributes,
$N = 1{,}000$ users, $s(A) = 0.05$, $s(B) = 0.04$,
$s(A,B) = 0.12$:

\begin{align*}
\lambda_{AB} &= \frac{0.12}{0.05 \times 0.04}
= \frac{0.12}{0.002} = 60 \\[4pt]
\Rraw &= 0.05 + 0.04 + 60 \times 0.05 \times 0.04 \\
&= 0.05 + 0.04 + 0.12 = 0.21 \\[4pt]
R_{\mathrm{max}}(2) &= 2 + \binom{2}{2} = 2 + 1 = 3 \\[4pt]
R_{\mathrm{norm}} &= \frac{0.21}{3} = 0.070
\end{align*}

$R_{\mathrm{norm}} = 0.070$ means 7.0\% of the maximum
possible de-anonymization risk. Despite the very high
interaction coefficient $\lambda_{AB} = 60$, the absolute
risk remains low because the individual uniqueness scores
are small ($s(A) = 0.05$, $s(B) = 0.04$). The
normalisation correctly captures that amplification of
a small base risk produces a small absolute risk.
This illustrates a key insight: a high $\lambda_{ij}$
is only dangerous when the underlying $s(a_i)$ values
are also substantial.

\subsection{Interpretation Guide}

Table~\ref{tab:r_norm_interp} provides a practical
interpretation guide for $R_{\mathrm{norm}}$ values
in the CBDC deployment context.

\begin{table}[H]
\centering
\caption{Interpretation of $R_{\mathrm{norm}}$ values}
\label{tab:r_norm_interp}
\begin{tabular}{lll}
\toprule
\textbf{$R_{\mathrm{norm}}$ range} &
\textbf{Risk level} &
\textbf{Interpretation} \\
\midrule
$0.000$--$0.050$ & Very low  &
Fewer than 5\% of maximum risk \\
$0.051$--$0.100$ & Low       &
5--10\% of maximum risk \\
$0.101$--$0.200$ & Moderate  &
10--20\% of maximum risk \\
$0.201$--$0.350$ & High      &
20--35\% of maximum risk \\
$0.351$--$1.000$ & Very high &
Above 35\% of maximum risk \\
\bottomrule
\end{tabular}
\end{table}

Plain Hash Token with $R_{\mathrm{norm}} = 0.317$ falls
in the High category. ZKP Token with
$R_{\mathrm{norm}} = 0.035$ falls in the Very Low
category. Ring Signature and Hybrid Scheme with
$R_{\mathrm{norm}} \approx 0.069$ fall in the Low
category. These classifications directly inform the
policy recommendations in Chapter~\ref{ch:results}.

% ================================================================
\section{Mitigation Modelling}
\label{sec:mitigation}
% ================================================================

\subsection{Motivation}

The baseline model measures $R_{\mathrm{norm}}$ under
the assumption that all metadata is observable at its
natural granularity. In practice, privacy-enhancing
technical controls can reduce the observability or
uniqueness of individual attributes. This section defines
how mitigation strategies are incorporated into the model.

\subsection{Multiplicative Mitigation Approach}

Each mitigation strategy is represented as a set of
\textit{multipliers} $\mu_i \in [0, 1]$ applied to the
individual uniqueness scores:

\begin{equation}
s_{\mathrm{mit}}(a_i) = \mu_i \times s(a_i)
\label{eq:mitigation}
\end{equation}

A multiplier of $\mu_i = 1.0$ means no mitigation ---
the attribute is observed at full precision. A multiplier
of $\mu_i = 0.0$ means perfect suppression --- the
attribute contributes zero uniqueness. Intermediate values
represent partial mitigation. The joint uniqueness under
mitigation is approximated as:

\begin{equation}
s_{\mathrm{mit}}(a_i, a_j)
= \mu_i \times \mu_j \times s(a_i, a_j)
\label{eq:mit_joint}
\end{equation}

This multiplicative approximation assumes that mitigation
on two attributes scales the joint uniqueness by the
product of their individual multipliers, which is exact
under independence and a conservative approximation
under positive correlation.

\subsection{Justification of the Multiplicative Approach}

The multiplicative approach is chosen for three reasons.

First, it is \textit{interpretable}. A multiplier of
$\mu_i = 0.20$ on the amount attribute means that
differential privacy noise reduces observable uniqueness
to 20\% of its unmitigated level. This maps naturally
to the practical effect of adding Gaussian or Laplace
noise to transaction amounts.

Second, it is \textit{conservative}. By applying the
same multiplier to the joint uniqueness as to the
individual scores, the model does not underestimate
residual risk after mitigation.

Third, it is \textit{computationally efficient}. The
mitigated scores can be computed from the original
dataset by simple scalar multiplication, without
re-generating data or re-computing all uniqueness
scores from scratch.

\subsection{Mitigation Strategies Evaluated}

Seven mitigation strategies are evaluated in
Chapter~\ref{ch:results}. Table~\ref{tab:mitigations}
shows the multiplier configuration of each strategy.

\begin{table}[H]
\centering
\caption{Mitigation strategies and attribute multipliers}
\label{tab:mitigations}
\begin{tabular}{lp{8cm}}
\toprule
\textbf{Strategy} & \textbf{Multipliers applied} \\
\midrule
No Mitigation &
All $\mu_i = 1.00$. Baseline --- no controls applied. \\

Differential Privacy on Amount &
$\mu_{\mathrm{amount}} = 0.20$, all others $= 1.00$.
Simulates Laplace or Gaussian noise injection reducing
amount uniqueness by 80\%. \\

Timestamp Suppression &
$\mu_{\mathrm{timestamp}} = 0.05$, all others $= 1.00$.
Simulates coarsening timestamp to hour of day or day
of week, reducing temporal uniqueness by 95\%. \\

Wallet Unlinkability &
$\mu_{\mathrm{wallet}} = 0.05$, all others $= 1.00$.
Simulates stealth address or nullifier-based wallet
generation reducing wallet uniqueness by 95\%. \\

Nepal Controls &
$\mu_{\mathrm{channel}} = 0.20$,
$\mu_{\mathrm{zone}} = 0.15$, all others $= 1.00$.
Simulates coarsening of Nepal-specific attributes to
fewer categories (digital/physical; valley/non-valley). \\

DP and Timestamp &
$\mu_{\mathrm{amount}} = 0.20$,
$\mu_{\mathrm{timestamp}} = 0.05$,
all others $= 1.00$.
Combined differential privacy on amount and
timestamp suppression. \\

Full Stack &
$\mu_{\mathrm{amount}} = 0.20$,
$\mu_{\mathrm{timestamp}} = 0.05$,
$\mu_{\mathrm{merchant}} = 0.50$,
$\mu_{\mathrm{wallet}} = 0.05$,
$\mu_{\mathrm{frequency}} = 0.30$,
$\mu_{\mathrm{channel}} = 0.20$,
$\mu_{\mathrm{zone}} = 0.15$.
All seven attributes mitigated simultaneously. \\
\bottomrule
\end{tabular}
\end{table}

The Full Stack strategy represents the maximum technically
feasible privacy enhancement within the constraints of
a functional CBDC system that still supports AML
monitoring and tax reporting obligations.


% ================================================================
\section{Hyperledger Fabric Context}
\label{sec:fabric_context}
% ================================================================

\subsection{Why Hyperledger Fabric}

The reference deployment platform for this thesis is
Hyperledger Fabric, a permissioned enterprise blockchain
framework maintained under the Linux Foundation's Hyperledger
project \citep{hyperledgerfabric2023}. Fabric is chosen
because it is the most widely referenced platform in the
CBDC implementation literature for token-based retail
deployments \citep{biscgide2024}, and because its
architecture directly determines which metadata fields are
observable on the ledger.

Unlike public permissionless blockchains, Fabric requires
every participant to be enrolled through a Membership Service
Provider (MSP) before joining the network. This permissioned
model makes Fabric suitable for CBDC deployment where KYC
and AML compliance require verified participant identities.
Androulaki et al. \citeyear{androulaki2018} describe Fabric's
core architecture: modular consensus, channel-based data
partitioning, and chaincode (smart contract) execution in
isolated containers.

\subsection{Observable Metadata on a Fabric CBDC Ledger}

Every transaction on a Fabric channel produces a transaction
envelope containing several metadata fields visible to all
channel members. In a CBDC chaincode, the following fields
are directly observable:

\begin{itemize}
    \item \textbf{Transaction amount} --- the token transfer
    value, recorded as a chaincode argument.

    \item \textbf{Timestamp} --- the block orderer timestamp,
    recorded in the block header with the precision of the
    ordering service clock.

    \item \textbf{Merchant (recipient)} --- the recipient
    organisation ID or wallet address, recorded as a chaincode
    argument.

    \item \textbf{Wallet (sender)} --- the sender's MSP
    identity or token address, recorded in the transaction
    proposal header.

    \item \textbf{Frequency} --- not a single-field value,
    but derivable by any channel participant who counts
    transactions from a given wallet address across the
    ledger history.

    \item \textbf{Payment channel} --- derivable from the
    chaincode ID, which differs between mobile wallet,
    QR code, NFC, and offline token chaincodes.

    \item \textbf{Transaction zone} --- derivable from the
    endorsing peer's organisation, whose geographic MSP
    configuration identifies its zone.
\end{itemize}

These seven fields map directly to the $m = 7$ attributes
in this model. The choice of attributes is therefore not
arbitrary --- it reflects what a passive observer with read
access to a Fabric CBDC channel can observe without any
cryptographic attack.

\subsection{Permissioned Identity and the Privacy Problem}

Fabric's permissioned identity model means NRB, as the
network operator, knows the real-world identity behind every
wallet address. This is by design: KYC and AML compliance
require it. However, this does not eliminate the privacy
problem addressed in this thesis. As argued in
Chapter~\ref{ch:introduction}, the metadata privacy problem
exists at a different level --- it concerns what other
channel participants, insiders with partial access, and
adversaries who obtain the transaction log without the
identity mapping can infer from observable metadata alone.

\subsection{Scope Boundary: Domestic Transactions Only}

This model is scoped to domestic CBDC transactions settled
within Nepal's Fabric network under NRB jurisdiction.
Cross-border remittance transactions involve separate
settlement rails and are excluded. The seven modelled
attributes are all observable within the domestic channel;
remittance corridor is not observable within this boundary
and is therefore excluded from the attribute set.

% ================================================================
\section{Attribute Set Design}
\label{sec:attribute_design}
% ================================================================

\subsection{The m=7 Attribute Set}

This thesis models seven metadata attributes: five standard
attributes present in any CBDC deployment, and two
Nepal-specific attributes grounded in Nepal's payment
infrastructure and geographic structure.
Table~\ref{tab:attribute_set} defines the complete set.

\begin{table}[H]
\centering
\caption{The $m = 7$ attribute set with value domains}
\label{tab:attribute_set}
\begin{tabular}{llll}
\toprule
\textbf{Attribute} & \textbf{Type} &
\textbf{Value Domain} & \textbf{Nepal?}\\
\midrule
amount    & Continuous  & $\mathbb{R}^+$ (scheme-dependent granularity) & No\\
timestamp & Discrete    & Integer (scheme-dependent precision) & No\\
merchant  & Categorical & Finite set of IDs or categories & No\\
wallet    & Identifier  & Scheme-dependent address pool & No\\
frequency & Discrete    & $\mathbb{Z}^+$ bounded by scheme & No\\
payment\_channel  & Categorical & $\{0,1,2,3\}$ (mobile/QR/NFC/offline) & Yes\\
transaction\_zone & Categorical & $\{0,1,2,3,4\}$ (five geographic zones) & Yes\\
\bottomrule
\end{tabular}
\end{table}

\subsection{Nepal-Specific Attribute: Payment Channel}

Payment channel takes four values reflecting Nepal's digital
payment infrastructure. The distribution is informed by the
dominance of eSewa and Khalti in Nepal's mobile payment
ecosystem:

\begin{itemize}
    \item 0 = Mobile wallet (60\%) --- eSewa and Khalti
    \item 1 = QR code (25\%) --- rapid adoption at merchants
    \item 2 = NFC (10\%) --- urban smartphone users
    \item 3 = Offline token (5\%) --- rural low-connectivity
\end{itemize}

NFC and offline token users form small, identifiable
subpopulations. At small $N$, these users may be uniquely
identifiable from channel alone.

\subsection{Nepal-Specific Attribute: Transaction Zone}

Transaction zone takes five values reflecting Nepal's
geographic population distribution:

\begin{itemize}
    \item 0 = Kathmandu valley (35\%) --- capital region
    \item 1 = Hill urban (20\%) --- Pokhara, Biratnagar
    \item 2 = Hill rural (25\%) --- largest rural segment
    \item 3 = Terai urban (12\%) --- southern plains cities
    \item 4 = Terai rural (8\%) --- southern plains rural
\end{itemize}

Hill rural and Terai rural users form the smallest
subpopulations and are consequently the most vulnerable to
re-identification in small communities with limited merchants.

\subsection{Excluded Attribute: Remittance Corridor}

Remittance corridor was considered as a third Nepal-specific
attribute but excluded on scope grounds. This thesis models
a domestic CBDC under NRB jurisdiction. Cross-border
remittances involve separate settlement rails, foreign
correspondent banking relationships, and regulatory
frameworks outside NRB's domestic CBDC mandate. Remittance
corridor is not an attribute observable within the domestic
CBDC transaction record and therefore falls outside the
defined system boundary.

\subsection{R\_max for m=7}

With $m = 7$ attributes, the theoretical maximum is:

\begin{equation}
R_{\mathrm{max}}(7) = 7 + \binom{7}{2}
= 7 + 21 = 28
\end{equation}

All $\Rnorm$ values in this thesis are normalised against
this ceiling of 28.

% ================================================================
\section{The Six Cryptographic Schemes}
\label{sec:six_schemes}
% ================================================================

Six cryptographic schemes are modelled, spanning the full
spectrum of realistic token-based CBDC privacy designs.
Table~\ref{tab:six_schemes} summarises each scheme.

\begin{table}[H]
\centering
\caption{Six cryptographic schemes and their privacy
mechanisms}
\label{tab:six_schemes}
\begin{tabular}{lp{4cm}p{4cm}}
\toprule
\textbf{Scheme} & \textbf{Key Mechanism} &
\textbf{What It Hides}\\
\midrule
Plain Hash Token &
SHA-256 hash of token serial number &
Nothing --- all metadata fully exposed\\

Blind Signature &
Chaum's blind signature protocol &
Wallet randomised; timestamp coarsened to hour\\

Hybrid Scheme &
Blind signature wallet + ring signature amounts &
Wallet randomised; amounts bucketed to NPR 100\\

Ring Signature &
Ring anonymity set over $k$ possible signers &
Wallet in larger pool; amounts bucketed; fewer merchant categories\\

Stealth Address &
One-time wallet address per transaction &
Wallet completely unlinkable across transactions; amounts remain visible\\

ZKP Token &
Zero-knowledge range proofs; nullifier wallet &
All attributes reduced to coarse bands; wallet unlinkable\\
\bottomrule
\end{tabular}
\end{table}

\textbf{Plain Hash Token} is the weakest baseline. All
transaction metadata is recorded at full precision. The
wallet address is sequential and permanently fixed to each
user. No metadata transformation occurs.

\textbf{Blind Signature} applies Chaum's
\citeyear{chaum1983} blinding protocol. The bank signs a
token without seeing its serial number, enabling wallet
randomisation. The timestamp is coarsened to the hour of
day (24 values). Amounts and merchant IDs remain at full
precision.

\textbf{Hybrid Scheme} models the digital euro's proposed
tiered privacy design, combining wallet blinding from blind
signatures with amount bucketing from ring signatures.
Amounts are rounded to NPR 100 increments. Merchant
categories are reduced from 200 to 100. This scheme fills
the gap between Blind Signature and Ring Signature on the
risk ladder.

\textbf{Ring Signature} applies the CryptoNote ring
signature construction \citep{vansaberhagen2013}. The
wallet is drawn from a pool 20 times the population size.
Amounts are bucketed to NPR 100 increments. Merchant
categories are reduced to 20 broad types. This provides
a meaningful anonymity set for the wallet attribute.

\textbf{Stealth Address} generates a mathematically unique,
fresh address for every transaction using the formula:
$\text{wallet}(u) = u \times 997 + \epsilon$ where
$\epsilon \sim \text{Uniform}(0, 100{,}000)$. No two
transactions by the same user share a wallet address with
any mathematical relationship between them. This breaks
the spending pattern attack entirely for the wallet
attribute, but amounts and timestamps remain visible.

\textbf{ZKP Token} applies zero-knowledge range proofs
\citep{groth2015} to reduce all attributes to coarse
disclosure bands. Amounts are restricted to five
denominations $\{0, 100, 500, 1000, 5000\}$ NPR.
Timestamps are reduced to day of week (7 values). Merchants
are categorised into 5 broad economic sectors. The wallet
is a nullifier drawn from a pool of size $N \times 100$,
providing transaction unlinkability with a large anonymity
set.

% ================================================================
\section{Synthetic Data Generation}
\label{sec:data_gen}
% ================================================================

\subsection{Justification for Synthetic Data}

This thesis uses synthetically generated transaction data
rather than real NRB transaction records. This choice is
justified on three grounds. First, no real NRB CBDC
transaction dataset exists, as Nepal has not yet deployed
a CBDC. Second, synthetic data allows precise control over
the distributions of each attribute under each scheme,
enabling a clean comparison of cryptographic design choices
without confounding factors from real usage patterns.
Third, the distributions chosen are grounded in Nepal's
observed payment infrastructure context, making the
synthetic data a principled approximation of realistic
deployment conditions.

\subsection{Distribution Design}

For each scheme, the value domain of each attribute is
set to reflect how that scheme transforms the underlying
transaction data. Table~\ref{tab:data_gen} describes the
generation logic for each scheme-attribute combination.

\begin{table}[H]
\centering
\caption{Data generation logic by scheme and attribute}
\label{tab:data_gen}
\begin{tabular}{lp{8.5cm}}
\toprule
\textbf{Scheme} & \textbf{Generation logic}\\
\midrule
Plain Hash Token &
Amount: $\text{Exponential}(500)$ rounded to 2dp.
Timestamp: Uniform integer in $[0, 10N]$.
Merchant: Uniform from $\{0,\ldots,199\}$.
Wallet: Sequential $\{0,\ldots,N-1\}$.
Frequency: Uniform in $[1, 99]$.\\

Blind Signature &
Amount: $\text{Exponential}(500)$ rounded to 2dp.
Timestamp: Uniform hour in $[0, 23]$.
Merchant: Uniform from $\{0,\ldots,199\}$.
Wallet: Uniform from pool of size $10N$.
Frequency: Uniform in $[1, 49]$.\\

Hybrid Scheme &
Amount: $\text{Exponential}(500)$ rounded to NPR 100.
Timestamp: Uniform hour in $[0, 23]$.
Merchant: Uniform from $\{0,\ldots,99\}$.
Wallet: Uniform from pool of size $10N$.
Frequency: Uniform in $[1, 34]$.\\

Ring Signature &
Amount: $\text{Exponential}(500)$ rounded to NPR 100.
Timestamp: Uniform hour in $[0, 23]$.
Merchant: Uniform from $\{0,\ldots,19\}$.
Wallet: Uniform from pool of size $20N$.
Frequency: Uniform in $[1, 19]$.\\

Stealth Address &
Amount: $\text{Exponential}(500)$ rounded to 2dp.
Timestamp: Uniform hour in $[0, 23]$.
Merchant: Uniform from $\{0,\ldots,19\}$.
Wallet: $u \times 997 + \epsilon$, $\epsilon \sim
\text{Uniform}(0, 100{,}000)$.
Frequency: Uniform in $[1, 14]$.\\

ZKP Token &
Amount: Choice from $\{0, 100, 500, 1000, 5000\}$.
Timestamp: Choice from $\{0,\ldots,6\}$ (day of week).
Merchant: Choice from $\{0,\ldots,4\}$ (5 sectors).
Wallet: Uniform from pool of size $100N$.
Frequency: Choice from $\{0,\ldots,4\}$ (5 bands).\\
\bottomrule
\end{tabular}
\end{table}

Nepal-specific attributes are generated identically across
all schemes because they are observable at the network level
regardless of which cryptographic scheme is deployed on the
token layer:

\begin{itemize}
    \item payment\_channel: drawn from
    $\{0,1,2,3\}$ with probabilities
    $[0.60, 0.25, 0.10, 0.05]$
    \item transaction\_zone: drawn from
    $\{0,1,2,3,4\}$ with probabilities
    $[0.35, 0.20, 0.25, 0.12, 0.08]$
\end{itemize}

\subsection{Reproducibility}

All analyses use NumPy's \texttt{default\_rng} with a fixed
seed of 42 for the baseline analysis. The seed sensitivity
analysis (Analysis 9) re-runs all computations with seeds
0 through 29 to validate that results are not artefacts of
any particular random draw.

% ================================================================
\section{Implementation}
\label{sec:implementation}
% ================================================================

\subsection{Pipeline Architecture}

The complete model is implemented in Python 3 using NumPy,
Pandas, Matplotlib, and Seaborn. The core pipeline function
executes the following steps for any given dataframe and
attribute list:

\begin{lstlisting}
def pipeline(df, attrs):
    m    = len(attrs)
    s    = {a: s_single(df[a].values) for a in attrs}
    lmat = {a: {} for a in attrs}
    for ai, aj in combinations(attrs, 2):
        sij          = s_joint(df[ai].values, df[aj].values)
        l            = lam(s[ai], s[aj], sij)
        lmat[ai][aj] = lmat[aj][ai] = l
    Rr = R_raw(s, lmat, attrs)
    Rn = Rr / R_max(m)
    return {"s": s, "lmat": lmat, "R_raw": Rr,
            "R_norm": Rn, "m": m, "R_max": R_max(m)}
\end{lstlisting}

This function: (1) computes $s(a_i)$ for all $m$ attributes;
(2) computes $\lambda_{ij}$ for all $\binom{m}{2}$ pairs;
(3) accumulates $\Rraw$ using Equation~\ref{eq:r_raw};
(4) normalises to produce $\Rnorm$.

\subsection{Software Environment}

\begin{table}[H]
\centering
\caption{Software dependencies}
\label{tab:software}
\begin{tabular}{ll}
\toprule
\textbf{Library} & \textbf{Purpose}\\
\midrule
Python 3.10+  & Core language\\
NumPy         & Random number generation, array operations\\
Pandas        & Value counting, dataframe operations\\
Matplotlib    & Chart generation\\
Seaborn       & Heatmap visualisation\\
itertools     & Combinations generation for pair iteration\\
\bottomrule
\end{tabular}
\end{table}

The complete source code is provided in
Annex~\ref{app:a_code}. All analyses in
Chapter~\ref{ch:results} are produced by running a single
Python file: \texttt{cbdc\_nepal\_final.py}.